% Generated draft: requires mathematical and source review.
\documentclass[11pt]{article}
\usepackage[margin=25mm]{geometry}
\usepackage{fontspec}
\setmainfont{Latin Modern Roman}
\usepackage{amsmath,amssymb,mathtools,unicode-math}
\setmathfont{Latin Modern Math}
\usepackage{enumitem,needspace,xurl,hyperref}
\setlength{\parindent}{0pt}
\setlength{\parskip}{5pt}
\setlength{\emergencystretch}{3em}
\newcommand{\N}{\mathbb{N}}
\newcommand{\Z}{\mathbb{Z}}
\newcommand{\Q}{\mathbb{Q}}
\newcommand{\R}{\mathbb{R}}
\newcommand{\C}{\mathbb{C}}
\newcommand{\F}{\mathbb{F}}
\newcommand{\A}{\mathbb{A}}
\newcommand{\PP}{\mathbb{P}}
\newcommand{\E}{\mathbb{E}}
\newcommand{\eps}{\varepsilon}
\newcommand{\dd}{\,\mathrm{d}}
\begin{document}
\section*{UnsolvedMath ID: 1200025}
\textbf{Source problem code:} catalogue entry
\subsection{Definitions and mathematical statement}
For a countable group $\Gamma$, write $\beta_1^{(2)}(\Gamma)$ for its first $L^2$ Betti number, defined using group von Neumann dimension of the first reduced $L^2$ homology of a free contractible $\Gamma$-complex. The basic finite-index multiplicativity property is
\[
\beta_1^{(2)}(H)=[\Gamma:H]\,\beta_1^{(2)}(\Gamma)
\qquad(H\le\Gamma,\ [\Gamma:H]<\infty).
\]
The source asks for a proof of such properties by free spanning forest methods, not whether the already named multiplicativity property is true.

To specify the usual forest setting, let $\Gamma$ be infinite and finitely generated, let $S$ be a finite symmetric generating set, and let $G=\operatorname{Cay}(\Gamma,S)$ be its locally finite connected Cayley graph, rooted at the identity $o$. The free uniform spanning forest (FUSF) is the weak limit of uniform spanning trees on finite connected subgraphs exhausting $G$, with free boundary conditions, meaning that boundary vertices are not identified. If $\mathcal F_G$ has this law, $\deg_{\mathcal F_G}(o)$ is the number of incident retained edges. The supplied summary attributes to Lyons the formula, in this infinite-group setting,
\[
\E[\deg_{\mathcal F_G}(o)]=2\bigl(1+\beta_1^{(2)}(\Gamma)\bigr).
\]
For a finite-index subgroup $H$ of index $k$ and a chosen finite generating set $T$ for $H$, a forest proof of multiplicativity could therefore establish, without using multiplicativity as an input,
\[
\E[\deg_{\mathcal F_{\operatorname{Cay}(H,T)}}(o)]-2
=k\bigl(\E[\deg_{\mathcal F_G}(o)]-2\bigr).
\]
This is a concrete possible target implied by the reported formula, not a supplied forest proof. The infinite-group condition matters: an uncorrected expected-degree formula of this form cannot be applied to finite groups. The original methodological question does not state its group class or which other ``basic properties'' are intended.

The supplied summary claims a contraction-based equivalent identity involving a graph $G/W$, resistance and cycle maxima, a finite contraction identity attributed to Foster's theorem, and forest-based proofs for free and amenable groups. However, $W$, the special generating sets, and the displayed identities labeled $(\ast)$ and $(\dagger)$ are absent. Its finite-approximation and mass-transport discussion is therefore a reported proposed approach, not a specified reduction that can be verified here.

The source also distinguishes FUSF from the free minimal spanning forest (FMSF), ordinarily obtained from independent continuous edge labels by deleting the maximum-labeled edge on every finite cycle. Its claimed inequality involving FMSF expected degree and group cost has no complete hypotheses or reference in the record and is not adopted as an exact theorem here. The summary ends mid-sentence; no missing conclusion is inferred.
\subsection{Short English statement}
Can spanning-forest arguments prove established identities for first $L^2$ Betti numbers, especially scaling by subgroup index? In the infinite finitely generated setting, the supplied FUSF expected-degree formula turns this into an expected-degree identity between subgroup and ambient Cayley graphs. The record reports proposed reductions and special cases but omits their actual constructions and proofs.
\subsection{Sources}
\begin{itemize}
\item Local source registry: \url{data/unsolvedmath/problems.json}, permanent UnsolvedMath ID $1200025$, source code catalogue entry, including the supplied PARTIAL-PROGRESS summary. Dataset attribution: Ulam AI and UnsolvedMath Contributors; original-source attribution and rights are retained.
\item Abert, \emph{Some questions} (2010), Question 25: \url{https://www.renyi.hu/~abert/questions.pdf}. The supplied metadata reports manual transcription and July 29, 2026 access and flags current status and rights for review.
\item The supplied summary attributes the FUSF expected-degree formula to Lyons, the finite contraction identity to Foster's theorem, and the free-group special case to Nielsen--Schreier. It provides no complete bibliographic citations or separate links for these facts.
\item The same summary claims an equivalence labeled $(\ast)$, a contraction identity labeled $(\dagger)$, special-case forest proofs, and a proposed FMSF--cost bound, but supplies neither those formulas in full nor the referenced working arguments. These claims are not independently verified in this draft.
\end{itemize}
\end{document}
